\chapter{État de l'art}

\section{État de l'art sur le domaine d'application}

L'estimation de la valeur d'un bien immobilier repose historiquement sur le \textbf{modèle
hédonique}, formalisé par \citet{rosen1974hedonic}. Ce modèle postule qu'un bien composite
comme un logement peut être décomposé en un ensemble de caractéristiques (surface, nombre de
pièces, localisation, qualité de construction, etc.), chacune contribuant implicitement au prix
final observé sur le marché. Cette approche, traditionnellement opérationnalisée par une
régression linéaire multiple, reste la référence en économie urbaine et sous-tend encore la
méthode dite « des comparables » utilisée par les agents immobiliers et notaires : le prix d'un
bien est estimé par analogie avec des transactions récentes de biens aux caractéristiques
proches.

Cette approche présente cependant des limites bien documentées : elle suppose des relations
linéaires entre caractéristiques et prix, peine à capturer les interactions complexes entre
variables (par exemple, l'effet de la surface sur le prix n'est pas le même selon le quartier),
et dépend fortement de la disponibilité de transactions comparables récentes, une condition
difficile à remplir dans un marché comme celui du Sénégal, où les données de transactions ne
sont pas centralisées.

L'essor de l'apprentissage automatique a permis l'émergence des \emph{Automated Valuation
Models} (AVM), qui remplacent ou complètent le modèle hédonique classique par des algorithmes
capables de capturer des relations non linéaires et des interactions complexes entre
caractéristiques, sans hypothèse a priori sur la forme fonctionnelle de la relation entre les
variables explicatives et le prix.

\section{État de l'art sur les travaux de recherche pertinents liés au sujet}

Plusieurs familles de modèles ont été appliquées à la prédiction de prix immobiliers dans la
littérature récente.

Les \textbf{méthodes ensemblistes à base d'arbres de décision} occupent une place centrale.
\citet{breiman2001randomforests} a introduit le \emph{Random Forest}, qui combine un grand
nombre d'arbres entraînés sur des sous-échantillons aléatoires des données et des variables,
réduisant la variance par rapport à un arbre unique. \citet{chen2016xgboost} a ensuite proposé
XGBoost, une implémentation optimisée du \emph{gradient boosting}, qui construit les arbres de
façon séquentielle pour corriger les erreurs des précédents. Appliqués à la prédiction de prix
immobiliers, \citet{truong2020housing} comparent régression linéaire, k plus proches voisins,
CART et méthodes ensemblistes, et montrent la supériorité systématique des approches
ensemblistes sur les modèles linéaires classiques. \citet{sharma2024optimal} comparent plus
spécifiquement régression linéaire multiple, SVR, Random Forest, MLP et XGBoost sur le jeu de
données Ames Housing (Iowa, États-Unis) et concluent à la supériorité de XGBoost, avec un
$R^2$ proche de 0,89, un résultat obtenu toutefois sur un jeu de données nettement plus riche
(79 variables descriptives) que celui mobilisé dans ce projet.

Les \textbf{réseaux de neurones profonds} ont également été explorés. \citet{yazdani2021mldl}
comparent modèles hédoniques classiques, méthodes de machine learning (Random Forest, k plus
proches voisins) et de deep learning pour la prédiction de prix immobiliers, et observent que
l'avantage des réseaux de neurones dépend fortement du volume de données disponible : sur des
jeux de données de taille modeste, les méthodes ensemblistes à base d'arbres restent souvent
compétitives, voire supérieures, aux architectures de deep learning.

Enfin, un travail directement pertinent pour le contexte de ce projet est celui de
\citet{nwankwo2023lagos}, qui appliquent plusieurs modèles de machine learning (dont un réseau
de neurones artificiel et XGBoost) à la prédiction des prix immobiliers à Lagos, au Nigeria, un
marché ouest-africain partageant avec le Sénégal des caractéristiques structurelles proches
(données peu centralisées, forte hétérogénéité des annonces en ligne). Les auteurs confirment
l'efficacité des approches ensemblistes et des réseaux de neurones dans ce contexte, tout en
soulignant le rôle déterminant des variables de localisation.

L'écosystème logiciel de ce projet s'appuie par ailleurs sur la bibliothèque scikit-learn
\citep{pedregosa2011scikitlearn}, largement utilisée dans la littérature pour l'implémentation
reproductible de pipelines de machine learning.

\section{Identification des zones de recherche manquantes liées au sujet}

Cette revue de la littérature met en évidence plusieurs zones encore peu explorées, que ce
projet cherche en partie à adresser :

\begin{itemize}
  \item \textbf{Absence de données et de travaux spécifiques au Sénégal.} Contrairement au
  Nigeria \citep{nwankwo2023lagos} ou aux marchés occidentaux disposant de jeux de données de
  référence (Ames Housing, Boston Housing), aucun jeu de données immobilier sénégalais n'est
  publiquement disponible, et aucun travail de recherche appliquant le machine learning à ce
  marché n'a été identifié au moment de la rédaction de ce document.
  \item \textbf{Comparaison systématique méthodes ensemblistes / deep learning sur données
  ouest-africaines restreintes.} La littérature suggère que l'avantage du deep learning sur les
  méthodes ensemblistes dépend du volume de données \citep{yazdani2021mldl}, mais ce compromis
  n'a pas été quantifié sur un marché immobilier africain à faible volume de données
  structurées.
  \item \textbf{Accessibilité des outils produits.} La grande majorité des travaux cités se
  limitent à une évaluation académique des modèles, sans restitution accessible à un utilisateur
  non technique, un manque que ce projet adresse via le développement d'une interface graphique
  interactive.
\end{itemize}

\section{Conclusion}

Le modèle hédonique reste le cadre théorique de référence pour l'estimation immobilière, mais
les approches d'apprentissage automatique, en particulier les méthodes ensemblistes à base
d'arbres, s'imposent progressivement comme des alternatives plus performantes, sous réserve
d'un volume de données suffisant. Le contexte sénégalais, marqué par l'absence de données
publiques structurées, constitue à la fois la principale difficulté méthodologique de ce projet
et sa principale zone de contribution : la constitution d'un premier jeu de données exploitable
et la comparaison empirique de plusieurs familles de modèles sur ce marché spécifique.
